\chapter{L6 -- Blocchi di Jordan, autovalori complessi e autovettori generalizzati}

\section{Esponenziale di un blocco di Jordan nilpotente}

Nella lezione precedente abbiamo visto che, se una matrice è simile a una forma di Jordan, allora il calcolo di $e^{At}$ si riduce al calcolo dell'esponenziale dei singoli blocchi di Jordan.  
Conviene quindi partire dal caso più semplice: il blocco di Jordan \emph{nilpotente} associato all'autovalore nullo.

\subsection{Blocco nilpotente elementare}

Indichiamo con
\[
J_0^{(p)}=
\begin{pmatrix}
0 & 1 & 0 & \cdots & 0\\
0 & 0 & 1 & \cdots & 0\\
\vdots & \vdots & \vdots & \ddots & \vdots\\
0 & 0 & 0 & \cdots & 1\\
0 & 0 & 0 & \cdots & 0
\end{pmatrix}\in\mathbb{R}^{p\times p}
\]
il blocco di Jordan nilpotente di ordine $p$.

Per esempio, per $p=3$:
\[
J_0^{(3)}=
\begin{pmatrix}
0 & 1 & 0\\
0 & 0 & 1\\
0 & 0 & 0
\end{pmatrix},
\qquad
\left(J_0^{(3)}\right)^2=
\begin{pmatrix}
0 & 0 & 1\\
0 & 0 & 0\\
0 & 0 & 0
\end{pmatrix},
\qquad
\left(J_0^{(3)}\right)^3=0.
\]

In generale:
\begin{equation}
\left(J_0^{(p)}\right)^p=0.
\label{eq:L6_nilpotent_power_zero}
\end{equation}

Questa proprietà dice che $J_0^{(p)}$ è \emph{nilpotente} di indice $p$.

\subsection{Esponenziale del blocco nilpotente}

Dalla definizione di esponenziale di matrice:
\[
e^{J_0^{(p)}t}
=
\sum_{k=0}^{\infty}\frac{\left(J_0^{(p)}\right)^k t^k}{k!}.
\]

Ma, grazie a \eqref{eq:L6_nilpotent_power_zero}, tutti i termini con $k\ge p$ sono nulli. Quindi la serie si tronca:
\begin{equation}
e^{J_0^{(p)}t}
=
I+\frac{J_0^{(p)}t}{1!}+\frac{\left(J_0^{(p)}\right)^2 t^2}{2!}+\cdots+\frac{\left(J_0^{(p)}\right)^{p-1} t^{p-1}}{(p-1)!}.
\label{eq:L6_exp_nilpotent_jordan}
\end{equation}

Per $p=3$ si ottiene
\[
e^{J_0^{(3)}t}
=
I+J_0^{(3)}t+\frac{\left(J_0^{(3)}\right)^2t^2}{2!}
=
\begin{pmatrix}
1 & t & \dfrac{t^2}{2}\\
0 & 1 & t\\
0 & 0 & 1
\end{pmatrix}.
\]

In generale, l'esponenziale del blocco nilpotente ha la forma
\begin{equation}
e^{J_0^{(p)}t}
=
\begin{pmatrix}
1 & t & \dfrac{t^2}{2!} & \cdots & \dfrac{t^{p-1}}{(p-1)!}\\
0 & 1 & t & \cdots & \dfrac{t^{p-2}}{(p-2)!}\\
0 & 0 & 1 & \ddots & \vdots\\
\vdots & \vdots & \ddots & \ddots & t\\
0 & 0 & \cdots & 0 & 1
\end{pmatrix}.
\label{eq:L6_exp_nilpotent_general}
\end{equation}

\paragraph{Interpretazione dinamica.}
Questo risultato è molto importante: un blocco nilpotente genera fattori polinomiali in $t$.  
Più il blocco è grande, più compaiono potenze alte di $t$ nella soluzione.

\section{Esponenziale di un blocco di Jordan generico}

Consideriamo ora un blocco di Jordan associato a un autovalore reale $d$:
\begin{equation}
J_d^{(p)}=
\begin{pmatrix}
d & 1 & 0 & \cdots & 0\\
0 & d & 1 & \cdots & 0\\
\vdots & \vdots & \vdots & \ddots & \vdots\\
0 & 0 & 0 & \cdots & 1\\
0 & 0 & 0 & \cdots & d
\end{pmatrix}.
\label{eq:L6_jordan_block_real}
\end{equation}

Si può scrivere come
\[
J_d^{(p)}=dI+J_0^{(p)}.
\]

Poiché $dI$ commuta con $J_0^{(p)}$, vale
\[
e^{J_d^{(p)}t}=e^{(dI+J_0^{(p)})t}=e^{dt}e^{J_0^{(p)}t}.
\]

Usando \eqref{eq:L6_exp_nilpotent_general}:
\begin{equation}
e^{J_d^{(p)}t}
=
e^{dt}
\begin{pmatrix}
1 & t & \dfrac{t^2}{2!} & \cdots & \dfrac{t^{p-1}}{(p-1)!}\\
0 & 1 & t & \cdots & \dfrac{t^{p-2}}{(p-2)!}\\
0 & 0 & 1 & \ddots & \vdots\\
\vdots & \vdots & \ddots & \ddots & t\\
0 & 0 & \cdots & 0 & 1
\end{pmatrix}.
\label{eq:L6_exp_jordan_real_block}
\end{equation}

\paragraph{Conseguenza qualitativa.}
L'andamento temporale di un blocco di Jordan reale non è governato soltanto da $e^{dt}$, ma da
\[
e^{dt},\qquad te^{dt},\qquad \frac{t^2}{2!}e^{dt},\ \dots
\]
fino al grado $p-1$.  
Quindi la dimensione del blocco conta, non solo il valore dell'autovalore.

\section{Potenze di un blocco di Jordan nel tempo discreto}

Nel tempo discreto la dinamica libera è
\[
x(k)=A^{k-k_0}x(k_0).
\]
Se $A$ è simile a una matrice di Jordan, allora bisogna capire come è fatto $J^k$.

Per un blocco reale
\[
J_d^{(p)}=dI+J_0^{(p)},
\]
vale il binomio matriciale
\begin{equation}
\left(J_d^{(p)}\right)^k
=
\sum_{i=0}^{p-1}\binom{k}{i}d^{\,k-i}\left(J_0^{(p)}\right)^i,
\label{eq:L6_power_jordan_block}
\end{equation}
dove la somma si ferma a $p-1$ perché $\left(J_0^{(p)}\right)^p=0$.

\paragraph{Caso $p=2$.}
Se
\[
J_d^{(2)}=
\begin{pmatrix}
d & 1\\
0 & d
\end{pmatrix},
\]
allora
\[
\left(J_d^{(2)}\right)^k
=
\begin{pmatrix}
d^k & k d^{k-1}\\
0 & d^k
\end{pmatrix}.
\]

\paragraph{Caso $p=3$.}
Se
\[
J_d^{(3)}=
\begin{pmatrix}
d & 1 & 0\\
0 & d & 1\\
0 & 0 & d
\end{pmatrix},
\]
allora
\[
\left(J_d^{(3)}\right)^k
=
\begin{pmatrix}
d^k & \binom{k}{1}d^{k-1} & \binom{k}{2}d^{k-2}\\
0 & d^k & \binom{k}{1}d^{k-1}\\
0 & 0 & d^k
\end{pmatrix}.
\]

\paragraph{Interpretazione dinamica nel discreto.}
Qui compaiono fattori del tipo
\[
d^k,\qquad kd^{k-1},\qquad \binom{k}{2}d^{k-2},\dots
\]
quindi non ci sono esponenziali continui, ma potenze e fattori polinomiali in $k$.

\paragraph{Conseguenza asintotica.}
Nel tempo discreto un blocco di Jordan associato a un autovalore reale $d$ converge a zero se
\[
|d|<1.
\]
Se invece $|d|>1$, diverge.  
Il caso $|d|=1$ è delicato: la presenza di blocchi di Jordan di ordine superiore a $1$ può generare crescita polinomiale.

\section{Autovalori complessi e necessità di una forma reale}

Fino ad ora abbiamo trattato il caso di autovalori reali.  
Se però $A\in\mathbb{R}^{n\times n}$ ha autovalori complessi, bisogna fare attenzione: lo stato del sistema è reale, quindi la descrizione finale dell'evoluzione deve restare reale.

\subsection{Coppie coniugate}

Se $A$ è reale e possiede un autovalore complesso
\[
\lambda=\sigma+j\omega,
\]
allora anche il suo coniugato
\[
\bar\lambda=\sigma-j\omega
\]
è autovalore di $A$.

Analogamente, se $v$ è autovettore associato a $\lambda$, allora $\bar v$ è autovettore associato a $\bar\lambda$.

Scriviamo
\[
v=q_r+jq_i,
\qquad
\bar v=q_r-jq_i,
\]
dove $q_r,q_i\in\mathbb{R}^n$ sono parte reale e parte immaginaria dell'autovettore.

\subsection{Caso semplice \texorpdfstring{$2\times 2$}{2x2}}

Supponiamo che $A\in\mathbb{R}^{2\times 2}$ abbia autovalori
\[
\lambda=\sigma+j\omega,
\qquad
\bar\lambda=\sigma-j\omega.
\]

Se
\[
Q=[\,v\ \bar v\,],
\qquad
D=
\begin{pmatrix}
\lambda & 0\\
0 & \bar\lambda
\end{pmatrix},
\]
allora
\[
AQ=QD.
\]

Il problema è che $Q$ è complessa. Per costruire una base reale si introduce la matrice
\[
E=\frac12
\begin{pmatrix}
1 & -j\\
1 & j
\end{pmatrix},
\]
così che
\[
QE=[\,q_r\ q_i\,]\in\mathbb{R}^{2\times 2}.
\]

Facendo i conti si ottiene
\[
E^{-1}DE
=
\begin{pmatrix}
\sigma & \omega\\
-\omega & \sigma
\end{pmatrix}
=:M.
\]

Quindi
\[
AQE=QDE=QE\,(E^{-1}DE)=QE\,M,
\]
e pertanto
\begin{equation}
A=(QE)\,M\,(QE)^{-1},
\label{eq:L6_real_reduction_complex}
\end{equation}
dove $QE$ è reale e
\begin{equation}
M=
\begin{pmatrix}
\sigma & \omega\\
-\omega & \sigma
\end{pmatrix}.
\label{eq:L6_real_block_complex_pair}
\end{equation}

\paragraph{Messaggio importante.}
Una coppia di autovalori complessi coniugati viene sostituita, in una base reale, con un blocco reale $2\times 2$.

\section{Esponenziale del blocco reale associato a una coppia complessa}

Consideriamo
\[
M=
\begin{pmatrix}
\sigma & \omega\\
-\omega & \sigma
\end{pmatrix}
=
\sigma I+
\begin{pmatrix}
0 & \omega\\
-\omega & 0
\end{pmatrix}.
\]

Poiché la parte diagonale commuta con la parte antisimmetrica, si ottiene
\begin{equation}
e^{Mt}
=
e^{\sigma t}
\begin{pmatrix}
\cos(\omega t) & \sin(\omega t)\\
-\sin(\omega t) & \cos(\omega t)
\end{pmatrix}.
\label{eq:L6_exp_real_complex_block}
\end{equation}

\paragraph{Interpretazione.}
La soluzione è il prodotto di:
\begin{itemize}
    \item un fattore esponenziale reale $e^{\sigma t}$, che governa crescita o decadimento;
    \item una rotazione/oscillazione con pulsazione $\omega$.
\end{itemize}

Quindi:
\begin{itemize}
    \item se $\sigma<0$, si hanno oscillazioni smorzate;
    \item se $\sigma=0$, si hanno oscillazioni persistenti;
    \item se $\sigma>0$, si hanno oscillazioni amplificate.
\end{itemize}

\subsection{Caso discreto}

Nel tempo discreto, se $\lambda=\sigma+j\omega$, è spesso utile passare alla forma polare
\[
\lambda=\rho(\cos\theta+j\sin\theta),
\qquad
\rho=\sqrt{\sigma^2+\omega^2}.
\]

Per il blocco reale
\[
M=
\begin{pmatrix}
\sigma & \omega\\
-\omega & \sigma
\end{pmatrix}
\]
vale
\begin{equation}
M^k
=
\rho^k
\begin{pmatrix}
\cos(k\theta) & \sin(k\theta)\\
-\sin(k\theta) & \cos(k\theta)
\end{pmatrix}.
\label{eq:L6_power_real_complex_block}
\end{equation}

\paragraph{Interpretazione discreta.}
Nel discreto:
\begin{itemize}
    \item $\rho$ regola l'ampiezza;
    \item $\theta$ regola la rotazione a ogni passo.
\end{itemize}

Quindi:
\begin{itemize}
    \item se $\rho<1$, la traiettoria tende a zero;
    \item se $\rho=1$, rimane limitata e puramente rotante;
    \item se $\rho>1$, diverge.
\end{itemize}

\section{Esempio con autovalori complessi}

Consideriamo
\[
A=
\begin{pmatrix}
3 & -1\\
2 & 1
\end{pmatrix}.
\]

Il polinomio caratteristico è
\[
p_A(\lambda)=\det(\lambda I-A)
=
(\lambda-3)(\lambda-1)+2
=
\lambda^2-4\lambda+5.
\]

Gli autovalori sono
\[
\lambda_{1,2}=2\pm j.
\]

Per $\lambda=2+j$ risolviamo
\[
(A-\lambda I)v=0.
\]

Si trova un autovettore complesso, per esempio
\[
v=
\begin{pmatrix}
1\\
1-j
\end{pmatrix}.
\]

Scrivendo
\[
v=q_r+jq_i
\]
si ha
\[
q_r=
\begin{pmatrix}
1\\
1
\end{pmatrix},
\qquad
q_i=
\begin{pmatrix}
0\\
-1
\end{pmatrix}.
\]

Quindi la matrice reale di cambio di base è
\[
Q_R=[\,q_r\ q_i\,]
=
\begin{pmatrix}
1 & 0\\
1 & -1
\end{pmatrix},
\]
e il blocco reale associato è
\[
M=
\begin{pmatrix}
2 & 1\\
-1 & 2
\end{pmatrix}.
\]

Ne segue che
\[
A=Q_R\,M\,Q_R^{-1},
\]
e quindi
\[
e^{At}=Q_R\,e^{Mt}\,Q_R^{-1}
=
Q_R\left[
e^{2t}
\begin{pmatrix}
\cos t & \sin t\\
-\sin t & \cos t
\end{pmatrix}
\right]Q_R^{-1}.
\]

\paragraph{Commento dinamico.}
Poiché la parte reale degli autovalori è $2>0$, il sistema ha andamento oscillatorio ma con ampiezza crescente: la dinamica libera esplode.

\section{Autovalore nullo e blocchi di Jordan}

Gli appunti insistono molto su un punto importante: l'autovalore da solo non basta, conta anche la dimensione del blocco di Jordan.

\subsection{Caso \texorpdfstring{$\lambda=0$}{lambda=0}}

\paragraph{Blocco semplice.}
Se
\[
J_0^{(1)}=(0),
\]
allora
\[
e^{J_0^{(1)}t}=1.
\]

Per la matrice nulla $2\times 2$:
\[
J=
\begin{pmatrix}
0 & 0\\
0 & 0
\end{pmatrix},
\qquad
e^{Jt}=I.
\]

Questa è una dinamica costante nel tempo: la soluzione non diverge ma neppure decade.

\paragraph{Blocco di ordine 2.}
Se
\[
J_0^{(2)}=
\begin{pmatrix}
0 & 1\\
0 & 0
\end{pmatrix},
\]
allora
\[
e^{J_0^{(2)}t}
=
\begin{pmatrix}
1 & t\\
0 & 1
\end{pmatrix}.
\]

Qui compare il termine $t$, che cresce senza limite.  
Quindi un autovalore nullo con blocco di Jordan non banale genera crescita polinomiale.

\paragraph{Messaggio.}
Nel continuo, per avere stabilità associata a un autovalore con parte reale nulla, è necessario che i blocchi di Jordan associati siano tutti di ordine $1$.

\subsection{Caso \texorpdfstring{$\lambda=-1$}{lambda=-1}}

Se
\[
J_{-1}^{(1)}=(-1),
\]
allora
\[
e^{J_{-1}^{(1)}t}=e^{-t},
\]
che tende a zero.

Se invece
\[
J_{-1}^{(2)}=
\begin{pmatrix}
-1 & 1\\
0 & -1
\end{pmatrix},
\]
allora
\[
e^{J_{-1}^{(2)}t}
=
e^{-t}
\begin{pmatrix}
1 & t\\
0 & 1
\end{pmatrix}
=
\begin{pmatrix}
e^{-t} & te^{-t}\\
0 & e^{-t}
\end{pmatrix}.
\]

Anche in questo caso la soluzione tende a zero, perché il decadimento esponenziale vince sul fattore polinomiale.

\section{Autovettori generalizzati}

Quando una matrice non è diagonalizzabile, gli autovettori non bastano per costruire una base.  
Per questo si introducono gli \emph{autovettori generalizzati}.

\subsection{Motivazione}

Consideriamo
\[
A=
\begin{pmatrix}
2 & 1 & 0\\
0 & 2 & 1\\
0 & 0 & 2
\end{pmatrix}
=
J_2^{(3)}.
\]

L'unico autovalore è $\lambda=2$, con molteplicità algebrica $3$.  
Calcoliamo
\[
A-2I=
\begin{pmatrix}
0 & 1 & 0\\
0 & 0 & 1\\
0 & 0 & 0
\end{pmatrix}.
\]

Il nucleo di $A-2I$ è formato dai vettori del tipo
\[
v=
\begin{pmatrix}
\alpha\\
0\\
0
\end{pmatrix},
\]
quindi
\[
\dim \ker(A-2I)=1.
\]

Dunque la molteplicità geometrica è $1$, minore della molteplicità algebrica $3$: la matrice non è diagonalizzabile.

\subsection{Definizione}

Un vettore $q$ si dice \emph{autovettore generalizzato di ordine $p$ associato all'autovalore $\lambda$} se
\begin{equation}
(A-\lambda I)^p q=0,
\qquad
(A-\lambda I)^{p-1}q\neq 0.
\label{eq:L6_generalized_eigenvector_def}
\end{equation}

Quindi:
\begin{itemize}
    \item per $p=1$ si ritrova il normale autovettore;
    \item per $p>1$ si ottengono i vettori necessari per completare le catene di Jordan.
\end{itemize}

\subsection{Autospazio generalizzato}

L'insieme
\[
\ker(A-\lambda I)^p
\]
è il sottospazio degli autovettori generalizzati di ordine al più $p$.

Si ha una catena crescente di sottospazi:
\begin{equation}
\ker(A-\lambda I)\subseteq \ker(A-\lambda I)^2 \subseteq \cdots \subseteq \ker(A-\lambda I)^p \subseteq \cdots
\label{eq:L6_chain_kernels}
\end{equation}

Questa catena si stabilizza quando si raggiunge la molteplicità algebrica dell'autovalore.

\section{Catene di autovettori generalizzati}

Per un blocco di Jordan di ordine $k$ associato all'autovalore $\lambda$, si costruisce una catena
\[
q^{(1)},q^{(2)},\dots,q^{(k)}
\]
tale che
\begin{equation}
(A-\lambda I)q^{(1)}=0,
\label{eq:L6_chain_1}
\end{equation}
\begin{equation}
(A-\lambda I)q^{(2)}=q^{(1)},
\label{eq:L6_chain_2}
\end{equation}
\[
\vdots
\]
\begin{equation}
(A-\lambda I)q^{(k)}=q^{(k-1)}.
\label{eq:L6_chain_k}
\end{equation}

Equivalentemente:
\[
Aq^{(1)}=\lambda q^{(1)},
\]
\[
Aq^{(2)}=q^{(1)}+\lambda q^{(2)},
\]
\[
\vdots
\]
\[
Aq^{(k)}=q^{(k-1)}+\lambda q^{(k)}.
\]

\begin{figure}[H]
\centering
\begin{tikzpicture}[>=Latex, node distance=1.65cm]
    \node[draw, rounded corners] (qk) {$q^{(k)}$};
    \node[draw, rounded corners, left=of qk] (qk1) {$q^{(k-1)}$};
    \node[left=of qk1] (dots) {$\cdots$};
    \node[draw, rounded corners, left=of dots] (q2) {$q^{(2)}$};
    \node[draw, rounded corners, left=of q2] (q1) {$q^{(1)}$};

    \draw[->, thick] (qk) -- node[above] {$A-\lambda I$} (qk1);
    \draw[->, thick] (qk1) -- (dots);
    \draw[->, thick] (dots) -- (q2);
    \draw[->, thick] (q2) -- node[above] {$A-\lambda I$} (q1);
    \draw[->, thick] (q1) -- ++(0,-1.2) node[below] {$0$};
\end{tikzpicture}
\caption{Catena di autovettori generalizzati associata a un blocco di Jordan.}
\end{figure}

\paragraph{Interpretazione.}
Una catena di ordine $k$ genera un blocco di Jordan di dimensione $k$.

\section{Significato della molteplicità geometrica}

Sia $\lambda_i$ un autovalore di $A$. Allora
\[
m_g(\lambda_i)=\dim\ker(A-\lambda_i I)
\]
è la molteplicità geometrica dell'autovalore.

\paragraph{Interpretazione strutturale.}
La molteplicità geometrica dice \emph{quanti blocchi di Jordan} sono associati a quell'autovalore.

In particolare:
\begin{itemize}
    \item se $m_g(\lambda_i)=m_a(\lambda_i)$, tutti i blocchi sono di ordine $1$ e la matrice è diagonalizzabile rispetto a quell'autovalore;
    \item se $m_g(\lambda_i)<m_a(\lambda_i)$, compaiono blocchi di Jordan di ordine superiore a $1$.
\end{itemize}

\section{Esempio di costruzione di una catena}

Riprendiamo
\[
A=
\begin{pmatrix}
2 & 1 & 0\\
0 & 2 & 1\\
0 & 0 & 2
\end{pmatrix}.
\]

Cerchiamo una catena per $\lambda=2$.

\subsection{Autovettore}

Serve
\[
(A-2I)q_1=0.
\]
Poiché
\[
A-2I=
\begin{pmatrix}
0 & 1 & 0\\
0 & 0 & 1\\
0 & 0 & 0
\end{pmatrix},
\]
si può scegliere
\[
q_1=
\begin{pmatrix}
1\\
0\\
0
\end{pmatrix}.
\]

\subsection{Autovettore generalizzato di ordine 2}

Cerchiamo $q_2$ tale che
\[
(A-2I)q_2=q_1.
\]
Una scelta naturale è
\[
q_2=
\begin{pmatrix}
0\\
1\\
0
\end{pmatrix},
\]
infatti
\[
(A-2I)q_2=
\begin{pmatrix}
1\\
0\\
0
\end{pmatrix}
=q_1.
\]

\subsection{Autovettore generalizzato di ordine 3}

Cerchiamo ora $q_3$ tale che
\[
(A-2I)q_3=q_2.
\]
Si può prendere
\[
q_3=
\begin{pmatrix}
0\\
0\\
1
\end{pmatrix},
\]
perché
\[
(A-2I)q_3=
\begin{pmatrix}
0\\
1\\
0
\end{pmatrix}
=q_2.
\]

\paragraph{Conclusione.}
La catena
\[
q_1,\ q_2,\ q_3
\]
costruisce esattamente il blocco $J_2^{(3)}$.

\section{Esempio con autovalore di molteplicità algebrica 3 e geometrica 2}

Consideriamo una matrice con autovalore $\lambda=5$ di molteplicità algebrica $3$ ma molteplicità geometrica $2$.  
Allora i possibili schemi di Jordan associati a $\lambda=5$ sono:
\begin{itemize}
    \item un blocco di ordine $3$;
    \item un blocco di ordine $2$ più un blocco di ordine $1$;
    \item tre blocchi di ordine $1$.
\end{itemize}

Se
\[
\dim\ker(A-5I)=2,
\]
allora i blocchi sono due, e quindi la struttura possibile è
\[
J=
\begin{pmatrix}
5 & 1 & 0\\
0 & 5 & 0\\
0 & 0 & 5
\end{pmatrix},
\]
cioè un blocco di ordine $2$ e uno di ordine $1$.

\paragraph{Procedura.}
\begin{enumerate}
    \item si trova una base di $\ker(A-5I)$, che fornisce gli autovettori di ordine $1$;
    \item si cerca poi un vettore in $\ker(A-5I)^2$ ma non in $\ker(A-5I)$;
    \item applicando $(A-5I)$ a questo vettore si ottiene uno degli autovettori del primo passo;
    \item questa costruisce una catena di lunghezza $2$.
\end{enumerate}

\section{Conteggio degli autovettori generalizzati}

Sia
\[
M_p:=\dim\ker(A-\lambda I)^p.
\]

Allora:
\begin{itemize}
    \item il numero di autovettori generalizzati di ordine almeno $1$ è $M_1$;
    \item il numero di autovettori generalizzati di ordine almeno $2$ è $M_2-M_1$;
    \item il numero di autovettori generalizzati di ordine almeno $3$ è $M_3-M_2$;
    \item e così via.
\end{itemize}

Più precisamente, il numero di catene di lunghezza almeno $p$ è
\[
M_p-M_{p-1}.
\]

\paragraph{Interpretazione.}
Questo permette di ricostruire la dimensione dei blocchi di Jordan semplicemente studiando la crescita delle dimensioni dei nuclei
\[
\ker(A-\lambda I),\ \ker(A-\lambda I)^2,\ \ker(A-\lambda I)^3,\dots
\]

\section{Schema pratico per costruire la forma di Jordan}

Se si vuole costruire una matrice $Q$ tale che
\[
Q^{-1}AQ=J,
\]
si procede così:

\begin{enumerate}
    \item si trovano gli autovalori di $A$ e le loro molteplicità algebriche;
    \item per ogni autovalore $\lambda_i$ si calcola
    \[
    \dim\ker(A-\lambda_i I),
    \quad
    \dim\ker(A-\lambda_i I)^2,
    \quad \dots
    \]
    fino a stabilizzazione;
    \item da queste dimensioni si deduce la struttura dei blocchi di Jordan;
    \item si costruiscono le catene di autovettori generalizzati;
    \item si mettono come colonne di $Q$ i vettori delle catene, ordinati in modo coerente con i blocchi di Jordan.
\end{enumerate}

\paragraph{Caso particolare: un solo autovalore e un solo blocco.}
Se la matrice ha un solo autovalore $\lambda$ con:
\[
m_a(\lambda)=n,\qquad m_g(\lambda)=1,
\]
allora esiste un unico blocco di Jordan di ordine $n$.  
In tal caso basta costruire una singola catena completa
\[
q_1,\ q_2,\ \dots,\ q_n.
\]

\section{Sintesi finale della lezione}

In questa lezione abbiamo visto che:

\begin{enumerate}
    \item l'esponenziale di un blocco nilpotente $J_0^{(p)}$ è un polinomio matriciale che si tronca al grado $p-1$;

    \item l'esponenziale di un blocco di Jordan reale è dato da
    \[
    e^{J_d^{(p)}t}=e^{dt}e^{J_0^{(p)}t},
    \]
    quindi compaiono fattori del tipo $t^k e^{dt}$;

    \item nel tempo discreto le potenze di un blocco di Jordan generano fattori del tipo
    \[
    \binom{k}{i}d^{k-i};
    \]

    \item per autovalori complessi coniugati, una coppia complessa si traduce in un blocco reale
    \[
    \begin{pmatrix}
    \sigma & \omega\\
    -\omega & \sigma
    \end{pmatrix},
    \]
    il cui esponenziale descrive oscillazioni smorzate, persistenti o amplificate;

    \item quando la matrice non è diagonalizzabile si introducono gli autovettori generalizzati;

    \item le catene di autovettori generalizzati permettono di costruire i blocchi di Jordan;

    \item la molteplicità geometrica di un autovalore coincide con il numero di blocchi di Jordan associati a quell'autovalore;

    \item studiando i nuclei di
    \[
    (A-\lambda I)^p
    \]
    si può ricostruire operativamente la forma di Jordan.
\end{enumerate}

Questa parte è decisiva per comprendere il comportamento dinamico dei sistemi lineari, soprattutto nei casi non diagonalizzabili e nei casi oscillatori con autovalori complessi.